\section{Flip classes of circular graphs}\label{sec:circular}

In this section we prove Theorem~\ref{thm:main} for circular graphs, by
reduction to the ordinary case.
The idea is to expand $\Kv(\theta)$ and $\Ev(\theta)$, for every admissible
orientation, as the same signed sum over \emph{balanced} sub-orientations, in
which the lengths of the edges are differences of a potential.
The coefficients turn out to be constant on flip classes, so the expansion can
be grouped by flip classes.
Finally, the orientations of a balanced flip class can be drawn on a line, where
they form a union of ordinary flip classes.

Let $\Gamma$ be a circular decorated unit-interval graph.
A strict edge is a corner, so it is never a short edge of a flip, and it is
present in every orientation.
Hence forgetting the strict markings, and including the strict edges in every
orientation, identifies the flip classes of $\Gamma$ with flip classes of the
undecorated graph, with the same graphs $\Dplus{\theta}$.
We therefore assume that $S=\emptyset$, and identify an orientation $\theta$ with
the edge set of $\Dplus{\theta}$.
Throughout, orientations are admissible, and for a set of edges $Y$ we write
$c(Y)$ for the number of connected components of the undirected graph on
$\setZ/n\setZ$ with edges $Y$, isolated vertices included.
Every edge $e$ has a \defin{length} $\ell(e)\in\{1,\dotsc,n-1\}$.

We write $\Kv(\theta)=\Kp_{\Pplus{\theta}}$, and $H_\theta$ for the whole map
$z\mapsto\widetilde{\hrvop}(\theta,z)$ on $\setZ$.
We work in the free abelian group on such maps, with basis elements $[H_\theta]$.
The linear map $[H]\mapsto\elementaryE_\lambda$, where $\lambda$ is the partition
of the fibre sizes of $H$ modulo $n$, sends $[H_\theta]$ to $\Ev(\theta)$.

\subsection{Balanced orientations}

\begin{definition}
A set $F$ of edges is \defin{balanced} if on each connected component of the
undirected graph $(\setZ/n\setZ,F)$ there is a function $p$ to $\setZ$ with
$\ell(e)=p(v)-p(u)$ for every edge $e$ from $u$ to $v$ in $F$.
\end{definition}

Equivalently, the lifted graph of $F$ is a disjoint union of translates of a
graph on finitely many integers.
Let $F\subseteq\theta$ be balanced, with such a potential $p$.
Contract the connected components of $F$, keeping the edges of $\theta\setminus F$
as individual edges, possibly loops or parallel edges, and give the edge $e$ from
$u$ to $v$ the \defin{gain} $\delta(e)=\ell(e)-(p(v)-p(u))$.
Changing $p$ by a constant on a component does not change the gain of a closed
walk.

We call a directed multigraph with gains \defin{negative} if every directed cycle
has negative gain, and \defin{cyclic} if every edge lies on a directed cycle.

\begin{definition}\label{def:coefficient}
For a balanced $F\subseteq\theta$, let
$\defin{\chi(\theta,F)}=(-1)^{c(F)-c(\theta)}$ if the contracted graph is negative
and cyclic.
Otherwise let $\chi(\theta,F)=0$.
For $F$ not balanced we put $\chi(\theta,F)=0$.
\end{definition}

If $\theta$ itself is balanced, then $\chi(\theta,F)=[F=\theta]$.
Indeed, restricting a potential of $\theta$ to $F\subsetneq\theta$ gives every edge
of $\theta\setminus F$ gain zero, so any directed cycle containing one of them,
a loop included, has gain zero.

\begin{example}\label{ex:deletion}
Let $\avec=(2,2,2)$, so that $\Gamma$ is the complete circular graph on three
vertices, and let $\theta$ consist of the edges $1\to2$, $1\to3$ and $3\to2$, of
lengths $1$, $2$ and $2$.
It is not balanced: the path $1\to3\to2$ has length $4$, while $1\to2$ has length
$1$.
Lemma~\ref{lem:deletion} below, applied to this cycle, gives
$\Kv(\theta)=\Kv(F)$ and $[H_\theta]=[H_F]$ for $F=\{1\to3,3\to2\}$.
The set $F$ is balanced; its lift is the path $1\to3\to5$, so all three
vertices reach $2$, and both values equal $\elementaryE_3$.
\end{example}

\subsection{The finite deletion identity}

\begin{lemma}\label{lem:deletion}
Let $\theta$ be admissible, and let $Z$ be a cycle of the undirected graph of
$\theta$, traversed so that its forward edges form a nonempty set $P$, its
backward edges form a set $N$, and
\[
 \sum_{e\in P}\ell(e)>\sum_{e\in N}\ell(e).
\]
Let $B=\theta\setminus N$.
Then
\[
 \sum_{J\subseteq N}(-1)^{|J|}\Kv(B\cup J)=0
 \qquad\text{and}\qquad
 \sum_{J\subseteq N}(-1)^{|J|}[H_{B\cup J}]=0.
\]
\end{lemma}

\begin{proof}
Let $\kappa$ be a coloring which increases along the edges of $B$.
The values of $\kappa$ telescope around $Z$, so
$\sum_{P}(\kappa(v)-\kappa(u))=\sum_{N}(\kappa(v)-\kappa(u))$, the sums ranging over
the edges from $u$ to $v$.
The left side is positive, so some edge of $N$ increases.
Toggling the first such edge, in a fixed order of $N$, is a sign-reversing
involution on the sets $J$ for which $\kappa$ increases along $B\cup J$, which
proves the first identity.

For the second, fix $J$ and let $H=H_{B\cup J}$.
The function $t(z)=H(z)-z$ is periodic, so it is a function on residues.
For an edge from $u$ to $v$ in $B\cup J$, its lift ends at $u+\ell$ and
$H(u)\geq H(u+\ell)$, so $t(u)-t(v)\geq\ell$.
Telescoping $t$ around $Z$ gives
\[
 \sum_{N}(t(u)-t(v))=\sum_{P}(t(u)-t(v))\geq\sum_P\ell>\sum_N\ell.
\]
Hence some edge of $N$ satisfies $t(u)-t(v)>\ell$, that is, $H(u)>H(u+\ell)$ for its
lifts.
Adding this edge does not change $H$, by Lemma~\ref{lem:ineffective}(i).
Removing it does not change $H$ either: a path through a lift of this edge ends
below $H(u+\ell)<H(u)$ and is not maximizing.
Toggling the first edge of $N$ with this property is therefore a sign-reversing
involution which preserves $H$.
\end{proof}

If $\theta$ is not balanced, extend a potential on a spanning forest of $\theta$.
An edge outside the forest violating the potential lies on a cycle of the forest
with nonzero signed length, and we traverse it in the direction in which the
forward edges are longer.
Then $N\neq\emptyset$, since otherwise $\theta$ would contain a directed cycle.
Hence every set $B\cup J$ with $J\subsetneq N$ in Lemma~\ref{lem:deletion} has fewer
edges than $\theta$.

\subsection{Expansion into balanced orientations}

An \defin{ordered partition} of the vertex set of a directed multigraph $G$ is
a sequence $\omega=(B_1,\dotsc,B_k)$ of disjoint nonempty blocks covering it; it is
\defin{compatible} with $G$ if every edge goes from some $B_i$ to some $B_j$ with
$i\leq j$.

\begin{lemma}\label{lem:orderedPartitions}
For every finite directed multigraph $G$,
\[
 \sum_{\omega}(-1)^{k(\omega)}
 =(-1)^{c(G)}\,[\text{every edge of } G \text{ lies on a directed cycle}],
\]
where $\omega$ ranges over the ordered partitions compatible with $G$, and
$k(\omega)$ is the number of blocks.
\end{lemma}

\begin{proof}
A strongly connected component lies in one block of a compatible ordered
partition, so we may contract the strong components and remove the loops.
The edges of $G$ all lie on directed cycles exactly when the resulting acyclic
graph $Q$ has no edges.
We show by induction that the sum for $Q$ is $(-1)^{|Q|}$ if $Q$ has no edges and
$0$ otherwise.
The last block $U$ of a compatible ordered partition is a nonempty set closed
under taking heads of edges, so the sum for $Q$ is minus the sum over such $U$ of
the sums for $Q\setminus U$.
By induction only the sets $U$ containing every head of an edge contribute, and
the sum becomes $-(-1)^{|Q|}\sum_U(-1)^{|U|}$ over the nonempty sets $U$
containing all heads.
If $Q$ has an edge, it has a source which is not a head, and the sum vanishes;
otherwise it equals $-1$.
\end{proof}

\begin{proposition}\label{prop:balancedExpansion}
For every admissible orientation $\theta$,
\[
 \Kv(\theta)=\sum_{F\subseteq\theta}\chi(\theta,F)\,\Kv(F)
 \qquad\text{and}\qquad
 [H_\theta]=\sum_{F\subseteq\theta}\chi(\theta,F)\,[H_F].
\]
In particular $\Ev(\theta)=\sum_F\chi(\theta,F)\,\Ev(F)$.
\end{proposition}

\begin{proof}
We proceed by induction on $|\theta|$.
If $\theta$ is balanced, then $\chi(\theta,F)=[F=\theta]$.
Otherwise choose a cycle as after Lemma~\ref{lem:deletion}.
That lemma expresses both values of $\theta$ through the sets $B\cup J$ with
$J\subsetneq N$, and by induction it suffices to prove
\begin{equation}\label{eq:coefficientCancellation}
 \sum_{J\subseteq N}(-1)^{|J|}\chi(B\cup J,F)=0
 \qquad\text{for all } F\subseteq\theta.
\end{equation}

Fix $F$; we may assume that $F$ is balanced, with a potential $p$.
Contract the components of $F$ and let $Q$ be the resulting multigraph with gains.
Its edges are the edges of $B\setminus F$, which we call \defin{required}, and the
edges of $N'=N\setminus F$, which are \defin{optional}.
Only the sets $J$ containing $F\cap N$ contribute.
Writing $Q_J$ for the required edges together with the optional edges in $J$,
and using $c(B\cup J)=c(Q_J)$, the identity~\eqref{eq:coefficientCancellation}
becomes, up to a common sign,
\begin{equation}\label{eq:gainSum}
 \sum_{J\subseteq N'}(-1)^{|J|+c(Q_J)}\,[Q_J \text{ is negative and cyclic}]=0.
\end{equation}
The cycle $Z$ projects to a closed walk in $Q$ which follows the required edges
forwards and the optional edges backwards, and since the potential
telescopes, its gain satisfies
\begin{equation}\label{eq:walkGain}
 \sum_{\text{required}}\delta-\sum_{\text{optional}}\delta
 =\sum_P\ell-\sum_N\ell>0,
\end{equation}
the sums on the left ranging over the edges of the walk.

By Lemma~\ref{lem:orderedPartitions}, the left side of~\eqref{eq:gainSum} is
$\sum_\omega(-1)^{k(\omega)}$ times the sum of $(-1)^{|J|}$ over the $J$ such that
$\omega$ is compatible with $Q_J$ and $Q_J$ is negative.
Fix $\omega$.
If an optional edge compatible with $\omega$ goes strictly upwards between blocks,
it lies on no cycle of a graph compatible with $\omega$, and toggling it is a
sign-reversing involution.
Otherwise every compatible optional edge stays in one block.
Along the closed walk, required edges do not decrease the block index,
compatible optional edges traversed backwards keep it, and incompatible optional
edges traversed backwards increase it strictly.
The total change is zero, so the walk lies in a single block, and it remains to
show that
\begin{equation}\label{eq:negativeSum}
 \sum_{J}(-1)^{|J|}\,[Q_J \text{ is negative}]=0,
\end{equation}
the sum ranging over the sets $J$ of compatible optional edges.

Let $M$ exceed both the number of vertices and the number of optional edges of
$Q$, and put $\gamma(e)=M\delta(e)+1$.
A simple cycle of length $r$ has $\gamma$-gain $M g+r$, where $g$ is its gain, so
$Q_J$ is negative if and only if no cycle has positive $\gamma$-gain.
In that case let $\tau(y)$ be the maximum of $0$ and the $\gamma$-gains of the
directed walks ending at $y$.
It is the least nonnegative function with $\tau(v)-\tau(u)\geq\gamma(e)$ for all edges
$e$ from $u$ to $v$ of $Q_J$.
We call an optional edge \defin{slack} if $\tau(v)-\tau(u)>\gamma(e)$.
Summing $\tau(v)-\tau(u)-\gamma(e)$ around the closed walk,
by~\eqref{eq:walkGain} the optional edges exceed the required ones by
$M(\sum\delta-\sum\delta)+(\#\text{required}-\#\text{optional})>0$,
so some optional edge of the walk is slack.
Adding a slack edge preserves the inequalities, so $\tau$ does not change, and
removing it does not change $\tau$ either, since a walk through it is not
maximizing.
Toggling the first slack edge is therefore a sign-reversing involution, which
proves~\eqref{eq:negativeSum} and the proposition.
\end{proof}

\subsection{Constancy on flip classes}

Let $C$ be a flip class of $\Gamma$, and for an admissible orientation $\eta$ put
\[
 a_C(\eta)=\sum_{\alpha\in C}\chi(\alpha,\eta).
\]

\begin{lemma}\label{lem:constancy}
The function $a_C$ is constant on every flip class.
\end{lemma}

\begin{proof}
Let $\theta$ and $\eta$ differ by a flip which exchanges $ab$ for $bc$ across
$ac$, in a window $a<b<c<a+n$ of the lift.
Let $\alpha\in C$ with $\chi(\alpha,\theta)\neq0$, so $\theta\subseteq\alpha$.
Choose a potential $p$ of $\theta$, contract, and let $\tau$ be the function of the
proof of Proposition~\ref{prop:balancedExpansion} for the gains of
$\alpha\setminus\theta$.
Then $z(v)=p(v)+\tau(v)/M$ satisfies $z(v)-z(u)=\ell(e)$ for the edges of $\theta$ and
$z(v)-z(u)>\ell(e)$ for the edges of $\alpha\setminus\theta$.
The lengths satisfy $\ell(ac)=\ell(ab)+\ell(bc)$, so $bc\notin\alpha$.
Let $\alpha'=\alpha-ab+bc$.

A lifted path in $\alpha$ between the endpoints of $ab$ has total length
$\ell(ab)$ and total increase of $z$ equal to $\ell(ab)$, so all its edges satisfy
the equality, that is, they lie in $\theta$.
Since $ab$ and $ac$ are covers for $\theta$, they are covers for $\alpha$, and
similarly $bc$ and $ac$ are covers for $\alpha'$, using $\eta$.
The function $z$ increases along every edge of $\alpha'$, so $\alpha'$ is admissible,
and $\alpha'$ differs from $\alpha$ by a flip; hence $\alpha'\in C$.
The orientations $\theta$ and $\eta$ have the same components and the same
potential, and $\alpha'\setminus\eta=\alpha\setminus\theta$, so
$\chi(\alpha',\eta)=\chi(\alpha,\theta)$.
The same construction from $\eta$ inverts $\alpha\mapsto\alpha'$, which proves
$a_C(\theta)=a_C(\eta)$.
\end{proof}

Summing Proposition~\ref{prop:balancedExpansion} over $C$ and grouping by flip
classes $B$, we obtain
\begin{equation}\label{eq:grouping}
 \Kv(C)=\sum_B a_C(B)\,\Kv(B),\qquad \Ev(C)=\sum_B a_C(B)\,\Ev(B),
\end{equation}
where $a_C(B)$ denotes the common value on $B$, and $a_C(B)=0$ unless $B$
contains a balanced orientation.

\subsection{Balanced classes are ordinary}

\begin{proposition}\label{prop:balancedClass}
Let $B$ be a flip class containing a balanced orientation.
Then $\Kv(B)=\Ev(B)$.
\end{proposition}

\begin{proof}
Let $h$ be a potential of a balanced orientation in $B$.
On each component, $h(v)-v$ is constant modulo $n$, so changing $h$ by constants
we may assume that $h(v)\equiv v\pmod n$ and that $h$ is injective.
The relation $\ell(ac)=\ell(ab)+\ell(bc)$ shows that $h$ is a potential for every
orientation in $B$.
Order the vertices by the values of $h$, and let $\Gamma'$ be the ordinary graph on
these ranks in which a pair is a natural edge if the corresponding lifts
$h(u)<h(v)$ form a lifted natural edge.
It is a unit-interval graph since the lifted natural edges are closed under
subintervals.

An orientation $\xi$ of $\Gamma'$ maps to the circular orientation consisting of the
residues of its lifted edges.
A lifted path starting at $h(u)$ only visits the chosen lifts $h(w)$, since
$h$ is injective modulo $n$.
Hence the image is admissible, its covers correspond to those of
$\Pplus{\xi}$, its colorings correspond to those of $\xi$, and its highest
reachable lifts are the values $h(\hrv{\xi}{u})$.
So the map preserves $\Kv$ and $\Ev$, and maps ordinary flips to circular flips;
the window has length less than $n$ since the long edge is natural.
Every orientation in $B$ is the image of some $\xi$, since $h$ is a potential for
it.
Hence the set of $\xi$ whose image lies in $B$ is closed under ordinary flips,
and is a union of ordinary flip classes, on which the ordinary case of
Theorem~\ref{thm:main} gives $\Kv=\Ev$.
\end{proof}

\begin{proof}[Proof of Theorem~\ref{thm:main} for circular graphs]
By~\eqref{eq:grouping} and Proposition~\ref{prop:balancedClass},
$\Kv(C)=\sum_Ba_C(B)\Kv(B)=\sum_Ba_C(B)\Ev(B)=\Ev(C)$.
\end{proof}
